\begin{afterword}
\summaryhead

The post asks whether a nonbeliever's awe at real things is a substitute for religion. Wanting to fly was never irrational; only the magic potions were. Likewise, if a brain scan showed the author's response to a space shuttle launch resembling a Christian's response to a nativity scene, that would not make space travel a religion. Calling it one is ``theomorphism,'' the assumption that every nonbeliever has a hole that wants filling.

Some atheists do have such a hole, the author concedes, and their imitations, ``hymns to the nonexistence of God,'' are always bad, because they are written only because churches have hymns. The test for any post-religious song or ritual is whether it would make sense if religion had never existed. Weddings pass; lectures on God's nonexistence fail; tears at a shuttle launch pass. Nor should anyone avoid things merely because they resemble religion. A God who spoke plainly would become boringly real, and a real God would remove the uncertainty that fervor makes up for. Awe at a real shuttle has no such weakness. The post ends: ``humanity actually exists.''

\responsehead

The post's main argument is sound: a feeling shared with religion does not make its object a religion, and the airplane analogy shows why. The problems are in the claims around that argument.

The verdict on atheist hymns is universal and unchecked. They are ``invariably awful'' and will ``without exception, suck.'' No example is named, so the reader cannot judge. The explanation is a claim about the motives of every such writer: that they write only because churches have hymns. The post adds that genuine religious art ``is not an imitation of anything.'' The hymn tradition the post praises also imitates: Isaac Watts titled a psalter \textsc{Psalms of David: Imitated in the Language of the New Testament}. So imitation cannot by itself explain a bad hymn, and the claim rests on the unnamed writers' motives.

The post also makes claims about believers in general. Religious fervor, it says, compensates for inner uncertainty, and belief shared by everyone would end the specialness of being one of the elect. Both are stated as fact without evidence. Experiments on compensatory conviction give the first some general support, but they concern social issues and group bias, not religious belief.

The post uses ``humanism'' in a broad sense, for secular awe at what humans do. It does not discuss organized humanism, whose first manifesto, in 1933, presented humanism as a new religion.

In short: a sound point that shared feelings do not make shared categories, surrounded by a universal verdict on atheist hymns with no example, and unsupported claims about believers' inner lives.
\end{afterword}
